% TOP_PROBLEM: {"id": 1119, "title": "Catalan-Mersenne Conjecture", "category_id": 1, "category": {"id": 1, "name": "number_theory", "display_name": "Number Theory", "description": "Properties of integers, prime numbers, Diophantine equations.", "slug": "number-theory", "order_index": 1, "created_at": "2026-07-31T15:26:25.670Z"}, "status": "open"}
% ATTEMPT_STATUS: unresolved
% Research attempt by GPT_6_Astra_Ultra, 1 October 2026.
% Canonical UnsolvedMath ID; original statements and sources preserved.
% FIRST_POSTED: 2026-10-01
% Imported from: unsolvedmath_additional_500.tex; UnsolvedMath ID 1119
% !TeX program = lualatex
% Compile twice: lualatex 1119.tex
% Self-contained: no external bibliography, images, or \input files.
% Numeric section labels and visible numbers are original UnsolvedMath IDs.
\documentclass[11pt,a4paper]{article}
\usepackage[margin=24mm,headheight=14pt]{geometry}
\usepackage{fontspec}
\setmainfont{Latin Modern Roman}
\setsansfont{Latin Modern Sans}
\setmonofont{Latin Modern Mono}
\usepackage{amsmath,amssymb,mathtools}
\usepackage{unicode-math}
\setmathfont{Latin Modern Math}
\usepackage{microtype}
\usepackage{enumitem,xurl,needspace}
\usepackage[dvipsnames]{xcolor}
\usepackage{hyperref}
\hypersetup{unicode=true,colorlinks=true,linkcolor=MidnightBlue,urlcolor=MidnightBlue,
 pdftitle={UnsolvedMath Problem 1119},pdfauthor={GPT 6 Astra Ultra},bookmarksnumbered=true}
\usepackage{bookmark,tocloft,titlesec,fancyhdr}
\setlength{\cftsecnumwidth}{4.2em}
\cftsetpnumwidth{2.5em}
\cftsetrmarg{4em}
\titleformat{\section}{\Large\bfseries\raggedright}{\thesection}{0.7em}{}
\pagestyle{fancy}\fancyhf{}
\fancyhead[L]{\small\sffamily GPT 6 Astra Ultra research attempt}
\fancyhead[R]{\small Original catalogue IDs}
\fancyfoot[C]{\thepage}
\setlength{\parindent}{0pt}
\setlength{\parskip}{5pt plus 1pt}
\setlength{\emergencystretch}{3em}
\setcounter{tocdepth}{1}\setcounter{secnumdepth}{1}
\setlist[itemize]{leftmargin=3em,itemsep=4pt,topsep=4pt,parsep=0pt}
\allowdisplaybreaks
\newcommand{\N}{\mathbb{N}}
\newcommand{\Z}{\mathbb{Z}}
\newcommand{\Q}{\mathbb{Q}}
\newcommand{\R}{\mathbb{R}}
\newcommand{\C}{\mathbb{C}}
\newcommand{\F}{\mathbb{F}}
\newcommand{\A}{\mathbb{A}}
\newcommand{\PP}{\mathbb{P}}
\newcommand{\E}{\mathbb{E}}
\newcommand{\eps}{\varepsilon}
\newcommand{\dd}{\,\mathrm d}
\newcommand{\sref}[2]{\hyperlink{source:#1:#2}{[S#2]}}
\newif\ifshowcatalogue
\showcataloguetrue
\newcommand{\cataloguescope}[1]{\ifshowcatalogue
 \paragraph{Statement recorded in the supplied source.}
 #1\fi}
\begin{document}
% UnsolvedMath ID 1119; source catalogue code NT-016

\setcounter{section}{1118}

\section{The Catalan--Mersenne Iteration}\label{1119}

{\small\sffamily UnsolvedMath ID 1119 · Number Theory}

\subsection{Definitions and mathematical statement}

\paragraph{Shared notation.}Unless a section says otherwise, $\N=\{1,2,\ldots\}$,
$\Z$ is the ring of integers, $\Q,\R,\C$ are the rational, real, and complex
fields, $[N]=\{1,\ldots,\lfloor N\rfloor\}$, and $\log$ is the natural
logarithm. For real functions of a parameter tending to infinity,
$f\ll g$ and $f=O(g)$ mean $|f|\leq Cg$ eventually for some fixed $C$;
$f\gg g$ reverses this comparison; $f=o(g)$ means $f/g\to0$;
$f\sim g$ means $f/g\to1$; and $f\asymp g$ means both order bounds.
A subscript on $O$ or $\ll$ specifies allowed constant dependence.
The expression $n^{o(1)}$ in an upper bound means growth smaller than
$n^\epsilon$ for each fixed $\epsilon>0$. Local definitions govern any
exception, finite-field convention, or use of the same symbol for another
object. An asymptotic claim is not automatically an effective algorithm.



Set $C_0=2$ and recursively $C_{n+1}=2^{C_n}-1$ for integers $n\geq0$. A composite number is an integer greater than one with a proper divisor greater than one. The displayed target concerns every index $n>4$, using this exact zero-based convention; no different indexing of the initial term is substituted. \sref{1119}{1} \sref{1119}{2}

\paragraph{Statement recorded in the supplied source.}
Are all Catalan-Mersenne numbers $C_n$ composite for $n > 4$? Here $C_0 = 2$ and $C_{n+1} = 2^{C_n} - 1$. \sref{1119}{1}

\subsection{Short English statement}

After the specified early terms, must every number in the iterated Mersenne construction be composite?

\subsection{Research attempt}

\paragraph{The precise target and indexing.}
This attempt follows the catalogue's specific assertion: every $C_n$
with $n>4$ is composite, with $C_0=2$. It does not substitute another
historical formulation about the sequence. Direct recursion gives
\[
 C_0=2,\quad C_1=3,\quad C_2=7,\quad C_3=127,\quad
 C_4=P=2^{127}-1,
\]
where
\[
 P=170141183460469231731687303715884105727.
\]
The next term is $C_5=2^P-1$, which is stored symbolically throughout.
The Mersenne arithmetic and Lucas--Lehmer criterion documented by the
GIMPS project [E1] were inspected on 1 October 2026. No claim about a
current global factor-search record is needed in this attempt.

\paragraph{Compositeness propagates, making the target one question.}
For integers $a,b>1$,
\[
 2^{ab}-1=(2^a-1)(1+2^a+2^{2a}+\cdots+2^{(b-1)a}).
\]
Both factors are greater than one. Thus if $C_j$ is composite, so is
$C_{j+1}$, and then every subsequent term by induction. It follows that
the catalogue assertion is equivalent to the single statement
\[
                         C_5\text{ is composite}.
\]
One implication specializes to $n=5$; the other is the displayed
factorization iterated. Proving some later term composite is insufficient:
it leaves $C_5$ and possibly several further terms undecided. Conversely,
a proof that $C_5$ is prime would disprove this exact catalogue target,
even if a later term were composite. This distinction avoids an indexing
shift in what must actually be settled.

\paragraph{Certifying the exponent without constructing the next term.}
The Lucas--Lehmer theorem for prime exponent $p>2$ says that $2^p-1$
is prime exactly when the sequence
\[
 s_0=4,\qquad s_{j+1}=s_j^2-2\pmod{2^p-1}
\]
has $s_{p-2}=0$. Using this standard criterion [E1], the script executes
the $125$ iterations for $p=127$, obtaining residue zero. The exponent
$127$ is prime by trial division by $2,3,5,7,11$. Thus $P$ is prime.
The criterion itself is cited rather than reproved; the modular
arithmetic instance is actually executed. This certificate concerns
$C_4$, not $C_5$. Running the corresponding criterion on $C_5$ would
involve $P-2$ iterations on numbers with about $P$ bits, a wholly
different computation.

\paragraph{The exact shape of every possible prime divisor.}
Let $q$ be a prime divisor of $2^P-1$. It is odd, and the multiplicative
order of two modulo $q$ divides the prime $P$. That order cannot be one,
since no prime divides $2-1$. Hence it is exactly $P$, so $P\mid q-1$.
Since $P$ and $q$ are odd, there is an integer $k\ge1$ with
\[
                         q=2kP+1.
\]
In particular a factor is already at least $2P+1$, even though the
number being factored is vastly larger still. Furthermore
$2^{(q-1)/2}=2^{kP}\equiv1\pmod q$; Euler's criterion and the
supplementary law for the quadratic character of two imply
$q\equiv1$ or $7\pmod8$. Since $P\equiv7\pmod8$, this restricts
$k$ to $0$ or $1\pmod4$. The residue restriction is a standard Mersenne
factor sieve also described in [E1]; the first order argument already
suffices for the complete candidate form.

These conditions do not imply that any candidate divides $C_5$.
In the opposite direction, any proper integer divisor, even a composite
one, would certify compositeness if its divisibility were proved. A
factor search need not prove every candidate prime before trying its
modular remainder.

\paragraph{A feasible exact factor test.}
For a specified $q$, binary modular exponentiation evaluates
$2^P\bmod q$ using only the $127$-bit exponent $P$. A residue of one
would show $q\mid C_5$. Thus trying a moderate list of candidate
factors is feasible even though writing down $C_5$ is not.
The script tests all integers $q=2kP+1$ for $1\le k\le10000$, without
a probable-prime filter and without omitting either residue class.
None divides $C_5$. Because every prime factor has this form, this
also excludes every prime divisor $q\le20000P+1$.
This lower bound on a possible factor is conditional only on the cited
Lucas--Lehmer theorem and the executed exact arithmetic, not on a
probabilistic primality test for the candidate $q$.

There is no resulting lower bound comparable to $\sqrt{C_5}$.
To prove primality by trial division would require eliminating every
prime divisor through that square root, whereas our upper candidate
is linear in $P$ and the square root is exponential in $P$. Conversely,
a failure to find a factor supplies no evidence strong enough to assert
compositeness. A heuristic expectation that a randomly chosen large
integer is usually composite has no force for this individual number.

\paragraph{The circular route through an exponent factor.}
One might hope to use the factorization $2^{ab}-1$ at the first
unresolved step. Its exponent here is $P$, already certified prime,
so that mechanism is unavailable. It becomes effective only after a
composite term has been found. Invoking it to establish the first
composite term would therefore assume the very input it needs.
Likewise, a formal factorization over complex roots of unity does not
give nontrivial integer factors when the exponent is prime.

\paragraph{Reproducibility, remaining gap, and outcome.}
The exponent certificate and factor scan are stored in
\texttt{checks.py} and \texttt{results.json} at
\path{_Local/Erdos_problems/UnsolvedMath/attempt_audit_2026_10_01/number_theory/}.
The decisive unresolved step is a compositeness certificate for
$2^P-1$, for example one nontrivial divisor, or a primality certificate
that would refute the stated assertion. Neither was obtained.
\textbf{Outcome: UNRESOLVED.} The infinite tail was reduced exactly to
its first term, and a finite candidate-factor interval was excluded;
the reduction and the finite scan are not a solution of that term.

\subsection{Sources}

{\small

\begin{itemize}

\item[\hypertarget{source:1119:1}{S1}] ULAM.AI, \emph{UnsolvedMath}, supplied \texttt{problems.json}, record 1119 (NT-016), statement and background. The embedded literature-review date is 2026-08-17; that is the archive’s review date, not a fresh certification. Dataset landing page: \url{https://huggingface.co/datasets/ulamai/UnsolvedMath}.

\item[\hypertarget{source:1119:2}{S2}] Double Mersenne number overview. \url{https://en.wikipedia.org/wiki/Double_Mersenne_number}. Bibliographic information imported from the supplied record.


\item[E1] Great Internet Mersenne Prime Search,
\emph{Mathematics and Research Strategy}, sections on trial factoring
and Lucas--Lehmer primality testing.
\url{https://www.mersenne.org/various/math.php}.
Project documentation inspected 1 October 2026. Used for the standard
Lucas--Lehmer criterion and Mersenne factor restrictions, not for a
claim about a completed primality test of $C_5$.

\end{itemize}

}

\label{end:1119}
\end{document}
